\section{消融与机制分析}
\label{sec:ablation}

\subsection{局部时序定义与统计假设}
表~\ref{tab:temporal_order} 先给出 ComGenVid 上同网格 patch 时序阶数的最直接对照。D=2 作为主线设定，D=1、D=3 和 D=4 都明显退化，因此本方法把局部证据锁定为二阶差分，而不是继续堆叠更高阶导数。

\begin{table}[!t]
  \centering
  \caption{ComGenVid 上同网格 patch temporal derivative order 对照。D=2 为当前 patch 主线。}
  \label{tab:temporal_order}
  \begin{tabular}{lcccc}
    \toprule
    阶数 & AUC & AP & $\Delta$AUC & $\Delta$AP \\
    \midrule
    D=2 & 0.9273 & 0.9309 & 0.0000 & 0.0000 \\
    D=3 & 0.8853 & 0.8936 & -0.0420 & -0.0373 \\
    D=4 & 0.8866 & 0.8949 & -0.0407 & -0.0361 \\
    \bottomrule
  \end{tabular}
\end{table}


对于窗口内第 $t$ 帧、第 $i$ 个 patch token $P_{t,i}$，局部加速度为
\begin{equation}
 A_{t,i}=P_{t+2,i}-2P_{t+1,i}+P_{t,i}.
\end{equation}
196 个 patch 保持 14$\times$14 行优先网格顺序，CLS 与 register token 在计算前移除。$A_{t,i}$ 沿特征维 L2 归一化，再使用独立真实校准视频拟合的白化变换与高斯似然打分。窗口内对 patch 与时间位置取均值，并用真实窗口 CDF 转成 high-is-real percentile。这一 same-grid D2 定义检测局部表征轨迹的非平滑加速度，而不是跨位置做 hard/soft matching。

表~\ref{tab:patch_temporal_comprehensive} 把 D=1、D=2、D=3、D=4，以及 spatial-only 和 motion-hard/soft 放在同一协议下比较，进一步说明 D=2 不是任意高阶导数的替代品，而是局部时序分支的有效设定。

\begin{table*}[!t]
  \centering
  \caption{ComGenVid 局部 patch 证据的完整时序定义消融。所有行均使用真实视频百分位校准和 pairwise balanced AUC/AP；D=1 lag-1 补足一阶差分对照，D=2 same-grid 为当前主线。}
  \label{tab:patch_temporal_comprehensive}
  \resizebox{\textwidth}{!}{%
  \begin{tabular}{lllcccc}
    \toprule
    证据组 & 变体 & 局部时序定义 & AUC & AP & $\Delta$AUC vs D=2 & $\Delta$AP vs D=2 \\
    \midrule
    静态对照 & Spatial-only & patch spatial percentile & 0.8091 & 0.8194 & -0.1182 & -0.1115 \\
    一阶对照 & D=1 lag-1 & same\_grid\_lag1 & 0.8556 & 0.8617 & -0.0718 & -0.0692 \\
    一阶对照 & D=1 multi-lag & same\_grid\_multilag & 0.8546 & 0.8586 & -0.0728 & -0.0723 \\
    匹配扰动 & Motion-hard & motion hard control & 0.8430 & 0.8491 & -0.0843 & -0.0818 \\
    匹配扰动 & Motion-soft & motion soft control & 0.8181 & 0.8256 & -0.1092 & -0.1054 \\
    二阶主线 & \textbf{D=2 same-grid} & same\_grid\_second\_order & 0.9273 & 0.9309 & +0.0000 & +0.0000 \\
    高阶对照 & D=3 same-grid & same\_grid\_third\_order & 0.8853 & 0.8936 & -0.0420 & -0.0373 \\
    高阶对照 & D=4 same-grid & same\_grid\_fourth\_order & 0.8866 & 0.8949 & -0.0407 & -0.0361 \\
    \bottomrule
  \end{tabular}%
  }
\end{table*}


表~\ref{tab:patch_likelihood_assumption_audit} 则检验白化后局部表示是否足够接近高斯假设。second-order diff 的对角方差、非对角耦合和正态性压力测试都优于原始 patch token，因此该似然建模比直接在 patch token 上拟合更稳健。

\begin{table*}[!t]
  \centering
  \caption{Patch 似然统计假设诊断。实验使用 ComGenVid 真实视频 patch cache 抽样；AD 和 D'Agostino--Pearson (DP) 为坐标级正态性压力测试，covariance 和 direction cosine 用于检验白化后近似球形性。}
  \label{tab:patch_likelihood_assumption_audit}
  \resizebox{\textwidth}{!}{%
  \begin{tabular}{lccccccc}
    \toprule
    表示 & 行数 & 白化维度 & Cov diag mean & Offdiag RMS & AD pass@5\% & DP $p>0.01$ & Cos std / expected \\
    \midrule
    Patch token & 49920 & 1023 & 0.992 & 0.017 & 0.180 & 0.156 & 0.0346/0.0313 \\
    Region-pooled token & 24960 & 1024 & 1.047 & 0.010 & 0.039 & 0.023 & 0.0335/0.0312 \\
    Second-order diff & 49920 & 1024 & 0.975 & 0.009 & 0.883 & 0.906 & 0.0319/0.0312 \\
    \bottomrule
  \end{tabular}%
  }
\end{table*}


统一 K=1 核心消融进一步确认，PatchD2-only 的 Macro AP 为 0.8214；加入 0.1 PatchSpatial 后 Local 降到 0.8144，差值为 $-0.0070$，95\% 区间为 [$-0.0088,-0.0054$]。因此 Local 的有效信号主要来自 D2，PatchSpatial 只因冻结 clean 定义而保留。固定 beta 敏感性中 $\beta=0.05$ 与 $0.10$ 的 Macro AP 分别为 0.8725 和 0.8723，差异小于校准 seed 波动，不足以支持事后替换锁定值。

\begin{table*}[!t]
  \centering
  \caption{Local first- versus second-order temporal ablation under the locked LSTL protocol. All entries are AUC/real-positive AP; only the Local temporal derivative order changes.}
  \label{tab:local_d1_d2_ablation}
  \resizebox{\textwidth}{!}{%
  \begin{tabular}{lccccc}
    \toprule
    Variant & ComGenVid & VideoFeedback & GenVideo & Macro-3 & Macro $\Delta$AUC/$\Delta$AP \\
    \midrule
    Local first-order only & 0.8846/0.8848 & 0.7675/0.7712 & 0.8188/0.7952 & 0.8237/0.8171 & -- \\
    Local second-order only & 0.8907/0.8908 & 0.7783/0.7870 & 0.8279/0.8066 & 0.8323/0.8281 & +0.0086/+0.0111 \\
    Full model with first-order & 0.8966/0.9055 & 0.8563/0.8645 & 0.8635/0.8385 & 0.8721/0.8695 & -- \\
    LSTL & 0.8968/0.9064 & 0.8597/0.8689 & 0.8657/0.8416 & 0.8741/0.8723 & +0.0019/+0.0028 \\
    \bottomrule
  \end{tabular}%
  }
\end{table*}


\subsection{窗口数量与聚合}
历史同协议时间覆盖实验比较了 K=1、K=3、K=5 和全部非重叠窗口。K=3 分支均值达到 0.8694/0.8697；K=5 降至 0.8663/0.8672，全部非重叠窗口为 0.8712/0.8695，AP 未改善且推理量与时长耦合。K=3 的 Local bottom-2 和 mean/bottom-2 hybrid 分别为 0.8680/0.8684 与 0.8687/0.8691，均低于分支均值。最低窗口天然随 $K$ 增大而下降，即使真实重校准也更容易放大剪辑、模糊和解码扰动；因此 U0 只保留 K=3 mean。

按历史长度分组，K=3 在 10--20 秒和 20 秒以上视频上的 AP 分别提高约 0.0239 和 0.0177，而 2--4 秒组下降约 0.0024；这支持“额外覆盖主要帮助较长视频”。但 6--10 秒组下降约 0.0101，effective-$K=1$ 组也从重校准中获益，因此不能把全部增益归因于时长。

\begin{table}[!t]
  \centering
  \caption{既有 region 与窗口聚合敏感性。所有候选均使用 200-real 校准；MS 行继承历史对照的其余设置，MW 行继承历史 K=3 对照，因此只作开发集敏感性，不用于重选 U0。}
  \label{tab:u0_region_aggregation}
  \resizebox{\columnwidth}{!}{%
  \begin{tabular}{lll}
    \toprule
    方向 & 配置 & Macro AUC/AP \\
    \midrule
    Token region & Fine region 1 (MS1) & 0.8687/0.8695 \\
    Token region & Coarse region 2 (MS2) & 0.8671/0.8664 \\
    Token region & Fine+coarse (MS3) & 0.8685/0.8687 \\
    Token region & Regions 1+2+3 (MS4) & 0.8674/0.8672 \\
    Window aggregation & K=3 branch mean (MW2) & 0.8694/0.8697 \\
    Window aggregation & Local bottom-2 (MW3) & 0.8680/0.8684 \\
    Window aggregation & Local mean/bottom-2 hybrid (MW4) & 0.8687/0.8691 \\
    \bottomrule
  \end{tabular}}
\end{table}


表~\ref{tab:u0_region_aggregation} 说明 region 与聚合方式具有数据集依赖性，因此它们适合作为边界分析，而不是重选 U0 的自由度。

\begin{table*}[!t]
  \centering
  \caption{局部尺度和聚合方式敏感性。表中 $\Delta$AUC 相对同一数据集、同一超参数族内最佳设置计算，用于区分稳定设计和数据集相关边界。}
  \label{tab:hyperparameter_sensitivity_matrix}
  \resizebox{\textwidth}{!}{%
  \begin{tabular}{lllccc}
    \toprule
    超参数族 & 数据集 & 设置 & AUC & AP & $\Delta$AUC vs family best \\
    \midrule
    region & ComGenVid & region=1, mean & 0.9088 & 0.9075 & -0.0211 \\
     &  & region=2, mean & 0.9230 & 0.9227 & -0.0069 \\
     &  & region=3, mean & 0.9299 & 0.9288 & +0.0000 \\
    \midrule
     & GenVideo & region=1, mean & 0.7976 & 0.7985 & -0.0097 \\
     &  & region=2, mean & 0.8072 & 0.8092 & +0.0000 \\
     &  & region=3, mean & 0.7893 & 0.7909 & -0.0179 \\
    \midrule
     & VideoFeedback & region=1, mean & 0.8228 & 0.8302 & +0.0000 \\
     &  & region=2, mean & 0.7941 & 0.7946 & -0.0288 \\
     &  & region=3, mean & 0.7533 & 0.7551 & -0.0695 \\
    \midrule
    aggregation & GenVideo & region=2, bottom-k=0.20 & 0.7652 & 0.7571 & -0.0420 \\
     &  & region=2, bottom-k=0.50 & 0.7795 & 0.7746 & -0.0278 \\
     &  & region=2, mean & 0.8072 & 0.8092 & +0.0000 \\
    \midrule
     & VideoFeedback & region=1, bottom-k=0.20 & 0.7935 & 0.8067 & -0.0293 \\
     &  & region=1, bottom-k=0.50 & 0.8052 & 0.8172 & -0.0176 \\
     &  & region=1, mean & 0.8228 & 0.8302 & +0.0000 \\
    \midrule
    bottom-k & ComGenVid & region=3, bottom-k=0.10 & 0.9245 & 0.9285 & -0.0068 \\
     &  & region=3, bottom-k=0.15 & 0.9262 & 0.9300 & -0.0050 \\
     &  & region=3, bottom-k=0.20 & 0.9273 & 0.9309 & -0.0039 \\
     &  & region=3, bottom-k=0.25 & 0.9282 & 0.9316 & -0.0030 \\
     &  & region=3, bottom-k=0.30 & 0.9292 & 0.9323 & -0.0020 \\
     &  & region=3, bottom-k=0.35 & 0.9298 & 0.9327 & -0.0014 \\
     &  & region=3, bottom-k=0.40 & 0.9304 & 0.9330 & -0.0008 \\
     &  & region=3, bottom-k=0.50 & 0.9312 & 0.9334 & +0.0000 \\
    \bottomrule
  \end{tabular}%
  }
\end{table*}


表~\ref{tab:hyperparameter_sensitivity_matrix} 进一步显示 region、aggregation 和 bottom-k 的最优值并不共享同一个设置，因而只能据此描述局部空间尺度边界，不能把某个候选简单外推到全部数据集。

\subsection{共同运动残差}
我们在未池化 $[T-2,196,D]$ 加速度上计算空间均值共同运动比例，并测试先减共同运动、再独立拟合白化和 CDF 的 residual D2。Spatial-mean residual 将历史 K=3 对照的 Macro AUC/AP 从 0.8694/0.8697 降到 0.8626/0.8609，AP 差值为 $-0.0088$，95\% 区间 [$-0.0103,-0.0073$]，仅 4/20 个生成器不下降。真实与生成视频的共同运动比例方向在三个数据集间不一致；四个重点负迁移生成器也不存在统一高共同运动偏差。因此共同运动不是当前误差的通用解释，mean/median residual 均不进入最终方法。

\subsection{空间尺度与 DINO 层}
旧 region=2 已经是在计算 D2 前对 raw 14$\times$14 token 做 2$\times$2 average pooling，即真正的 7$\times$7 coarse D2，而不是 likelihood 后的分数池化。独立校准的 fine+coarse 融合 Macro AP 为 0.8687，未超过 fine-only；继续增加尺度没有证据基础。

中间层实验从同一次 DINO forward 提取 mid、late 和 final patch token，每层分别执行 residual 定义、L2 归一化、白化、似然和真实 CDF。最佳 late+final 融合相对进入阶段的基线只提高约 0.0002 AP，置信区间跨 0，且 VideoFeedback 与 GenVideo 各下降约 0.005。该增益不足以抵消多套校准参数和缓存复杂度，最终只使用 final layer。

\subsection{按生成器的机制分解}
按生成器展开可以更清楚地看到融合带来的边界。表~\ref{tab:per_generator_component_matrix} 显示，\method 虽在大多数生成器上优于单分支，但也存在少数 patch-only 更强或 global-only 更强的例外，说明“融合更好”只能作为总体结论，不能被夸大为逐生成器单调改善。

\begin{table*}[!t]
  \centering
  \caption{按生成器展开的 global-only、patch-only 与 \method 组件结果。表中同时报告 AUC/AP 和 \method 相对两个分支的 AUC 差值，用于显示融合收益与负迁移边界。}
  \label{tab:per_generator_component_matrix}
  \resizebox{\textwidth}{!}{%
  \begin{tabular}{llcccccccc}
    \toprule
    数据集 & 生成器 & Global AUC & Patch AUC & \method AUC & $\Delta_{\alpha-G}$ AUC & $\Delta_{\alpha-P}$ AUC & Global AP & Patch AP & \method AP \\
    \midrule
    ComGenVid & Sora & 0.8393 & 0.9171 & 0.9088 & +0.0695 & -0.0082 & 0.8421 & 0.9198 & 0.9089 \\
     & VEO3 & 0.8670 & 0.9376 & 0.9308 & +0.0638 & -0.0068 & 0.8686 & 0.9420 & 0.9333 \\
    \midrule
    VideoFeedback & AnimateDiff & 0.8462 & 0.8225 & 0.8647 & +0.0185 & +0.0422 & 0.8659 & 0.8223 & 0.8833 \\
     & Fast-SVD & 0.8772 & 0.8713 & 0.9014 & +0.0242 & +0.0302 & 0.8804 & 0.8965 & 0.9117 \\
     & LVDM & 0.8731 & 0.8630 & 0.9004 & +0.0272 & +0.0374 & 0.8998 & 0.8703 & 0.9132 \\
     & LaVie-base & 0.8336 & 0.7812 & 0.8326 & -0.0010 & +0.0514 & 0.8473 & 0.7785 & 0.8480 \\
     & ModelScope & 0.8329 & 0.8367 & 0.8661 & +0.0332 & +0.0294 & 0.8496 & 0.8394 & 0.8732 \\
     & Pika & 0.8319 & 0.8724 & 0.8863 & +0.0544 & +0.0138 & 0.8452 & 0.8925 & 0.9001 \\
     & SoRA-Clip & 0.8238 & 0.7887 & 0.8287 & +0.0049 & +0.0400 & 0.8253 & 0.7816 & 0.8328 \\
     & Text2Video-Zero & 0.8357 & 0.6636 & 0.7607 & -0.0749 & +0.0971 & 0.8385 & 0.6743 & 0.7783 \\
     & VideoCrafter2 & 0.9425 & 0.8494 & 0.9096 & -0.0329 & +0.0603 & 0.9493 & 0.8682 & 0.9265 \\
     & ZeroScope-576w & 0.7781 & 0.8796 & 0.8772 & +0.0991 & -0.0023 & 0.8127 & 0.8789 & 0.8832 \\
    \midrule
    GenVideo & Crafter & 0.8002 & 0.8012 & 0.8307 & +0.0305 & +0.0295 & 0.7800 & 0.7965 & 0.8002 \\
     & Gen2 & 0.8731 & 0.8697 & 0.9066 & +0.0335 & +0.0369 & 0.8840 & 0.8944 & 0.9229 \\
     & Lavie & 0.8473 & 0.8114 & 0.8593 & +0.0120 & +0.0479 & 0.8423 & 0.8097 & 0.8604 \\
     & ModelScope & 0.7763 & 0.8154 & 0.8343 & +0.0580 & +0.0190 & 0.7776 & 0.8200 & 0.8394 \\
     & MorphStudio & 0.8259 & 0.7768 & 0.8338 & +0.0078 & +0.0569 & 0.8361 & 0.7693 & 0.8392 \\
     & Show\_1 & 0.8178 & 0.8053 & 0.8433 & +0.0255 & +0.0380 & 0.8020 & 0.8024 & 0.8284 \\
     & Sora & 0.7050 & 0.8386 & 0.8316 & +0.1266 & -0.0070 & 0.7346 & 0.8738 & 0.8559 \\
     & WildScrape & 0.7232 & 0.7394 & 0.7594 & +0.0362 & +0.0200 & 0.6644 & 0.7077 & 0.6798 \\
    \bottomrule
  \end{tabular}%
  }
\end{table*}


表~\ref{tab:per_generator_bootstrap_delta} 进一步给出生成器级 paired bootstrap 区间。多数生成器的 $\Delta_{\alpha-G}$ 和 $\Delta_{\alpha-P}$ 都为正，但 Text2Video-Zero、VideoCrafter2、部分 Sora 与少数 Global-only 友好的样本仍然保留负迁移或跨零区间，因此本方法的增益是总体稳定而不是处处占优。

\begin{table*}[!t]
  \centering
  \caption{生成器级 paired bootstrap AUC 差值区间。$\Delta_{\alpha-G}$ 检验 \method 相对 global-only 的收益；$\Delta_{\alpha-P}$ 检验融合是否超过 patch-only。}
  \label{tab:per_generator_bootstrap_delta}
  \resizebox{\textwidth}{!}{%
  \begin{tabular}{llcccccc}
    \toprule
    数据集 & 生成器 & $\Delta_{\alpha-G}$ mean & 95\% CI & 结论 & $\Delta_{\alpha-P}$ mean & 95\% CI & 结论 \\
    \midrule
    ComGenVid & Sora & +0.0692 & [+0.0621, +0.0764] & 正向 & -0.0081 & [-0.0128, -0.0030] & 负向 \\
     & VEO3 & +0.0638 & [+0.0574, +0.0706] & 正向 & -0.0069 & [-0.0111, -0.0027] & 负向 \\
    \midrule
    GenVideo & Crafter & +0.0162 & [-0.0092, +0.0453] & 跨零 & +0.0391 & [+0.0193, +0.0588] & 正向 \\
     & Gen2 & +0.0302 & [+0.0207, +0.0398] & 正向 & +0.0356 & [+0.0294, +0.0416] & 正向 \\
     & Lavie & +0.0091 & [-0.0014, +0.0201] & 跨零 & +0.0469 & [+0.0400, +0.0539] & 正向 \\
     & ModelScope & +0.0612 & [+0.0471, +0.0779] & 正向 & +0.0157 & [+0.0050, +0.0259] & 正向 \\
     & MorphStudio & +0.0094 & [-0.0041, +0.0239] & 跨零 & +0.0540 & [+0.0444, +0.0642] & 正向 \\
     & Show\_1 & +0.0274 & [+0.0117, +0.0407] & 正向 & +0.0349 & [+0.0245, +0.0465] & 正向 \\
     & Sora & +0.0839 & [+0.0328, +0.1383] & 正向 & -0.0036 & [-0.0379, +0.0332] & 跨零 \\
     & WildScrape & +0.0388 & [+0.0223, +0.0549] & 正向 & +0.0177 & [+0.0043, +0.0291] & 正向 \\
    \midrule
    VideoFeedback & AnimateDiff & +0.0146 & [+0.0046, +0.0264] & 正向 & +0.0434 & [+0.0371, +0.0497] & 正向 \\
     & Fast-SVD & +0.0234 & [+0.0127, +0.0337] & 正向 & +0.0297 & [+0.0234, +0.0368] & 正向 \\
     & LVDM & +0.0263 & [+0.0199, +0.0333] & 正向 & +0.0378 & [+0.0342, +0.0415] & 正向 \\
     & LaVie-base & -0.0013 & [-0.0076, +0.0054] & 跨零 & +0.0517 & [+0.0475, +0.0563] & 正向 \\
     & ModelScope & +0.0362 & [+0.0309, +0.0417] & 正向 & +0.0278 & [+0.0248, +0.0309] & 正向 \\
     & Pika & +0.0582 & [+0.0527, +0.0635] & 正向 & +0.0116 & [+0.0080, +0.0151] & 正向 \\
     & SoRA-Clip & +0.0051 & [-0.0056, +0.0173] & 跨零 & +0.0410 & [+0.0331, +0.0493] & 正向 \\
     & Text2Video-Zero & -0.0709 & [-0.0767, -0.0649] & 负向 & +0.0946 & [+0.0904, +0.0991] & 正向 \\
     & VideoCrafter2 & -0.0304 & [-0.0354, -0.0253] & 负向 & +0.0579 & [+0.0544, +0.0614] & 正向 \\
     & ZeroScope-576w & +0.0924 & [+0.0836, +0.1009] & 正向 & -0.0009 & [-0.0052, +0.0033] & 跨零 \\
    \bottomrule
  \end{tabular}%
  }
\end{table*}


\subsection{联合融合与其他负结果}
固定线性融合之外，我们用五折 cross-fitting 生成真实校准视频的 OOF $(G,L)$，比较对称 Mahalanobis、单侧 lower-tail Mahalanobis 和 Gaussian copula。固定 $0.6G+0.4L$ 的历史对照 AP 为 0.8697；J2 和最佳 J3 分别只有 0.8437 与 0.8602。J2 虽降低 ComGenVid 和 VideoFeedback conflict subset 的阈值错误率，却使 GenVideo 冲突样本恶化，并同时降低 SoRA-Clip、Text2Video-Zero 和 VideoCrafter2。二维联合模型没有优于稳定的单调融合。

此前已完成的 lag-1/multi-lag、hard/soft matching、patch D3/D4、Global raw D3、三分支与 gate/clip/cap 也均未超过 same-grid D2 主线，因而不在 U0 上重复搜索。表~\ref{tab:u0_failed_directions} 汇总主要负结果。所有未通过预设增益、最差数据集和 bootstrap 门槛的组件均单独拒绝，不允许通过继续调权重或堆叠失败组件恢复。

\begin{table*}[!t]
  \centering
  \caption{预先限定的主要失败方向。历史结果均保留，但未通过准入的组件不进入锁定 U0。}
  \label{tab:u0_failed_directions}
  \resizebox{\textwidth}{!}{%
  \begin{tabular}{lcl}
    \toprule
    方向 & Macro-3 AUC/AP 或 AP 差值 & 决策依据 \\
    \midrule
    Spatial-mean residual D2 & $0.8626/0.8609$ ($\Delta$AP $-0.0088$) & 三数据集均下降，bootstrap CI 全负 \\
    Unified fine region1 (MS1) & $0.8687/0.8695$ ($\Delta$AP $-0.0001$) & 最接近基线但 CI 跨 0 \\
    Fine+coarse token fusion (MS3) & $0.8685/0.8687$ ($\Delta$AP $-0.0010$) & 未超过 fine token，停止更多尺度 \\
    Late+final layer fusion & $0.8721/0.8698$ ($\Delta$AP $+0.0002$) & CI 跨 0，VideoFeedback/GenVideo 各下降约 0.005 \\
    K=5 mean & $0.8663/0.8672$ & 额外窗口增加成本但低于 K=3 \\
    All non-overlap mean & $0.8712/0.8695$ & 未改善 AP，且时长相关计算不均衡 \\
    K=3 Local bottom-2/hybrid & $0.8680/0.8684$ / $0.8687/0.8691$ & lower-tail 聚合过度惩罚 GenVideo \\
    Joint Typicality J2/J3 & AP $0.8437/0.8602$ & 均低于固定 $0.6/0.4$ 的 $0.8697$ \\
    D3/D4、multi-lag、hard/soft matching & 均未超过同网格 D2 & 历史同协议消融已拒绝，不重复搜索 \\
    \bottomrule
  \end{tabular}
  }
\end{table*}


唯一允许的新性能候选是无超参数 OAS shrinkage covariance。表~\ref{tab:u0_oas_candidate} 显示它只改变 Macro AP $+0.000043$，95\% 区间为 [$-0.000089,+0.000174$]；五 seed AP 标准差还从 0.002210 增至 0.002245。该候选未达到 $+0.003$ 准入线，最终保留稳定有效秩 whitening，并关闭后续 covariance 搜索。

\begin{table}[!t]
  \centering
  \caption{无超参数 OAS covariance 候选。S0 为稳定 U0，S1 只替换 Local D2 covariance。}
  \label{tab:u0_oas_candidate}
  \begin{tabular}{lcc}
    \toprule
    数据集 & S0 AUC/AP & S1 AUC/AP \\
    \midrule
    ComGenVid & 0.8968/0.9064 & 0.8969/0.9065 \\
    VideoFeedback & 0.8597/0.8689 & 0.8599/0.8691 \\
    GenVideo & 0.8657/0.8416 & 0.8656/0.8415 \\
    Macro-3 & 0.8741/0.8723 & 0.8741/0.8723 \\
    \bottomrule
  \end{tabular}
\end{table}


\subsection{锁定合成注入与定位}
在 300 条真实视频上，我们预先锁定 K=3 窗口，并向中央 25\% patch 区域注入局部冻结、重复、闪烁、仿射抖动和时间平移；全帧冻结与合成切镜作为控制。表~\ref{tab:u0_injection} 中正的 score drop 表示注入后真实性降低。Local D2 只对局部冻结产生很弱的预期响应（最终分数 $+0.0023$），对闪烁、仿射抖动和时间平移的平均响应方向错误。全帧冻结使 Local 分数反而提高 0.4084，原因是冻结压低了二阶 token 运动；合成切镜也未形成预期的负真实性响应。

\begin{table}[!t]
  \centering
  \caption{300 条真实视频上的锁定合成注入。正值表示注入后真实性下降；负值表示错误方向。}
  \label{tab:u0_injection}
  \resizebox{\columnwidth}{!}{%
  \begin{tabular}{lrrrr}
    \toprule
    注入 & Global drop & Local drop & Final drop & AUPRC \\
    \midrule
    Local freeze & +0.0021 & +0.0027 & +0.0023 & 0.1558 \\
    Local repeat & +0.0031 & -0.0007 & +0.0016 & 0.0670 \\
    Local flicker & +0.0089 & -0.0153 & -0.0008 & 0.1032 \\
    Affine jitter & -0.0029 & -0.0106 & -0.0060 & 0.1030 \\
    Temporal shift & -0.0045 & -0.0023 & -0.0036 & 0.1245 \\
    Full-frame freeze & -0.0055 & -0.4084 & -0.1667 & 0.3424 \\
    Scene cut & -0.0182 & -0.0387 & -0.0264 & -- \\
    \bottomrule
  \end{tabular}}
\end{table}


Patch-time AUPRC 在局部冻结、重复、闪烁、仿射抖动和时间平移上分别为 0.1558、0.0670、0.1032、0.1030 和 0.1245。K=1 在 165/300 个视频中完全未覆盖注入帧，验证了时间覆盖不足；但 K=3 只增加扰动暴露概率，并不保证正确的异常方向。该实验不支持把 patch score map 解释成通用语义伪影定位器。

\subsection{最终复杂度选择}
最终 U0 仍只有 Global 与 Local 两个分支，且两者来自同一次冻结 DINO 前向。K=3 评测使用 787,690 个唯一帧，相对 K=1 的 342,736 个为 2.30 倍；K=5 和 all-window 分别为 3.21 倍和 4.02 倍。K=3 的 AP 增益有正向置信区间，而后两者没有，因此额外复杂度只保留到三个窗口。Residual、多尺度、中间层和 Joint 均不增加发布时推理路径。

\begin{table}[!t]
  \centering
  \caption{不同窗口覆盖的评测计算量。相对成本按唯一 DINO 帧数计算。}
  \label{tab:u0_compute_cost}
  \begin{tabular}{lrrr}
    \toprule
    配置 & 窗口数 & 唯一帧数 & 相对 K=1 \\
    \midrule
    K=1 & 21,421 & 342,736 & 1.00$\times$ \\
    K=3 & 56,812 & 787,690 & 2.30$\times$ \\
    K=5 & 92,098 & 1,100,359 & 3.21$\times$ \\
    All non-overlap & 86,096 & 1,377,536 & 4.02$\times$ \\
    \bottomrule
  \end{tabular}
\end{table}

